
\documentclass[12pt]{article}
\usepackage{amsmath}
\usepackage{amsfonts}
\usepackage{amssymb}
\usepackage{geometry}
\geometry{a4paper, margin=0.8in}
\pagenumbering{gobble}

\title{Homework solutions for 02YELMA \\ Chapter: Electrodynamics}
\begin{document}
\date{}
\maketitle
\vspace{-0.8in}

\begin{enumerate}

\item The magnetic field produced by a very long straight wire carrying current \(I\) is

\[
B(r)=\frac{\mu_0 I}{2\pi r}
\]

where \(r\) is the perpendicular distance from the wire.

Assume the square loop has side length \(a\), and its nearer side is at distance \(s\) from the wire.

\begin{itemize}
\item[(a)] Magnetic flux through the loop

Take a thin strip of width \(dr\) and height \(a\):

\[
dA=a\,dr
\]

so

\[
d\Phi = B(r)\,dA
\]

thus

\[
d\Phi
=
\frac{\mu_0 I}{2\pi r}\,a\,dr
\]

Integrate from \(r=s\) to \(r=s+a\):

\[
\Phi
=
\int_s^{s+a}
\frac{\mu_0 I a}{2\pi r}\,dr
\]

Therefore,

\[
\boxed{
\Phi
=
\frac{\mu_0 I a}{2\pi}
\ln\!\left(\frac{s+a}{s}\right)
}
\]

\item[(b)] Induced emf when the loop is pulled away

As the loop moves away, \(s=s(t)\) changes with speed

\[
\frac{ds}{dt}=v
\]

Using Faraday’s law,

\[
\mathcal{E}
=
-\frac{d\Phi}{dt}
\]

Differentiate:

\[
\mathcal{E}
=
-
\frac{\mu_0 I a}{2\pi}
\frac{d}{dt}
\ln\!\left(\frac{s+a}{s}\right)
\]

which gives

\[
\boxed{
\mathcal{E}
=
\frac{\mu_0 I a v}{2\pi}
\left(
\frac{1}{s}
-
\frac{1}{s+a}
\right)
}
\]

As the loop moves away, the magnetic flux decreases. By Lenz’s law, the induced current tries to maintain the original magnetic field, so the induced current is

\[
\boxed{\text{counterclockwise.}}
\]

\end{itemize}

\item Inside the solenoid,

\[
\vec B(t)=B_0\cos(\omega t)\,\hat z
\]

The loop radius is

\[
r=\frac{a}{2}
\]

so its area is

\[
A=\pi\left(\frac a2\right)^2
=
\frac{\pi a^2}{4}
\]

The magnetic flux is

\[
\Phi(t)=BA
=
B_0\cos(\omega t)\,
\frac{\pi a^2}{4}
\]

The induced emf is

\[
\mathcal E
=
-\frac{d\Phi}{dt}
=
\frac{\pi a^2 B_0\omega}{4}
\sin(\omega t)
\]

Using Ohm’s law,

\[
I=\frac{\mathcal E}{R}
\]

we obtain

\[
\boxed{
I(t)
=
\frac{\pi a^2 B_0\omega}{4R}
\sin(\omega t)
}
\]

\item The magnetic field inside a long solenoid is

\[
B=\mu_0 n I
\]

The flux through one turn is

\[
\Phi_1
=
BA
=
\mu_0 n I \,\pi R^2
\]

The total number of turns in length \(l\) is

\[
N=nl
\]

Thus the total flux linkage is

\[
N\Phi_1
=
nl\,
\mu_0 n I \pi R^2
\]

Using the definition of inductance,

\[
L=\frac{N\Phi_1}{I}
\]

we get

\[
L
=
\mu_0 n^2 \pi R^2 l
\]

Hence, the self-inductance per unit length is

\[
\boxed{
\frac{L}{l}
=
\mu_0 n^2 \pi R^2
}
\]

\item Electric field of a time-dependent solenoid

The magnetic field inside the solenoid is

\[
B(t)=\mu_0 n I(t)
\]

Using Faraday’s law,

\[
\oint \vec E\cdot d\vec l
=
-
\frac{d\Phi_B}{dt}
\]

By symmetry,

\[
\vec E = E_\phi(s)\,\hat\phi
\]

and therefore

\[
\oint \vec E\cdot d\vec l
=
E(2\pi s)
\]

\begin{itemize}
\item[(a)] Inside the solenoid (\(s<a\))

The magnetic flux is

\[
\Phi_B
=
B\pi s^2
\]

so

\[
E(2\pi s)
=
-
\pi s^2
\frac{dB}{dt}
\]

which gives

\[
\boxed{
E(s)
=
-\frac{s}{2}
\frac{dB}{dt}
=
-\frac{\mu_0 n s}{2}
\frac{dI}{dt}
}
\]

The direction is

\[
\boxed{\hat\phi}
\]

\item[(b)] Outside the solenoid (\(s>a\))

Only the field inside contributes:

\[
\Phi_B
=
B\pi a^2
\]

thus

\[
E(2\pi s)
=
-
\pi a^2
\frac{dB}{dt}
\]

so

\[
\boxed{
E(s)
=
-
\frac{a^2}{2s}
\frac{dB}{dt}
=
-
\frac{\mu_0 n a^2}{2s}
\frac{dI}{dt}
}
\]

Again, the direction is azimuthal (\(\hat\phi\)).
\end{itemize}

\item The applied voltage is

\[
V(t)=V_0\cos(\omega t)
\]

The electric field between the plates is

\[
E(t)=\frac{V(t)}{d}
=
\frac{V_0}{d}\cos(\omega t)
\]

\begin{itemize}
\item[(a)] Displacement current density

Using

\[
\vec J_d
=
\epsilon_0
\frac{\partial \vec E}{\partial t}
\]

we obtain

\[
\boxed{
\vec J_d(t)
=
-
\epsilon_0
\frac{V_0\omega}{d}
\sin(\omega t)\,\hat z
}
\]

\item[(b)] Magnetic field

Using Ampère-Maxwell law,

\[
\oint \vec B\cdot d\vec l
=
\mu_0\epsilon_0
\frac{d\Phi_E}{dt}
\]

and cylindrical symmetry,

\[
B(2\pi r)
=
\mu_0\epsilon_0
\frac{d}{dt}
(E \cdot \text{area})
\]

For \(r<R\), the enclosed area is

\[
\pi r^2
\]

thus

\[
B(2\pi r)
=
\mu_0\epsilon_0
\pi r^2
\frac{dE}{dt}
\]

which gives

\[
\boxed{
B(r,t)
=
-
\frac{\mu_0\epsilon_0 V_0\omega}{2d}
\,r\,
\sin(\omega t)
}
\]

For \(r>R\), the full plate area contributes:

\[
\pi R^2
\]

so

\[
\boxed{
B(r,t)
=
-
\frac{\mu_0\epsilon_0 V_0\omega}{2d}
\frac{R^2}{r}
\sin(\omega t)
}
\]

The magnetic field direction is

\[
\boxed{\hat\phi}
\]
\end{itemize}

\end{enumerate}


\end{document}
